\documentclass[letterpaper,10pt]{article}
\usepackage[left=2cm,top=3cm,right=2cm,bottom=4cm]{geometry} 
\usepackage[utf8]{inputenc}
\usepackage[activeacute,spanish]{babel}
\usepackage{multirow}	% Para poder unir filas en las tablas
\usepackage{colortbl}	% Para colorear tablas
\usepackage{latexsym,amssymb,oldlfont,amsmath}
\usepackage{graphicx,graphics}
\usepackage{multicol}
\usepackage{color}
\usepackage{listings}
\usepackage{lastpage}
\usepackage{enumitem}
\usepackage{mathtools}



%\usepackage{fontspec}
%\setmainfont{Times New Roman}
\newcommand{\R}{{\mathbb R}}
\newcommand{\si}{{\mathcal S}^n}
\renewcommand{\theequation}{\thesection.\arabic{equation}} % Cambia formato salida numeracion equation
\makeatletter\@addtoreset {equation}{section} \makeatother % Tambien se usa
\renewcommand{\baselinestretch}{1.3}
\usepackage{amssymb}
\usepackage{hyperref}
\definecolor{midnightblue}{rgb}{0.1,0.1,0.44}


\newcommand{\dis}{\displaystyle}
\renewcommand{\theenumii}{\arabic{enumii}}

\def\Z{\mathbb{Z}}
\def\N{\mathbb{N}}
\def\R{\mathbb{R}}
\def\M{\mathcal{M}_n(\mathbb{R})}
\def\sin{\operatorname{sen}}
\def\e{\operatorname{e}}
\DeclareMathOperator{\arcsec}{arcsec}
\DeclareMathOperator{\arccot}{arccot}
\DeclareMathOperator{\arccsc}{arccsc}
\providecommand{\abs}[1]{\lvert#1\rvert}


\usepackage{etoolbox,fancyhdr,xcolor}
\fancyhf{} % sets both header and footer to nothing
\renewcommand{\headrulewidth}{0pt}
%\renewcommand{\footrulewidth}{0.5pt}
\pagestyle{fancy}
\headheight=60pt %para cambiar el tamaño del encabezadox
\fancyhead[L] %la "L" indica a la izquierda
{
\begin{minipage}{6cm}
\includegraphics[width=0.9\textwidth]{logo_udp.eps}
\end{minipage}
}

%\fancyhead[C]
%{
%\textbf{Cálculo II - Guía 1}
%}


%\fancyhead[R] %la "L" indica a la izquierda
%{

%}
%\begin{minipage}{4cm}
%\small %tamaño de letra pequeño
%{
%\textsf
%{
%Universidad de Chile\\
%...}}
%\end{minipage}}



%\fancyfoot[C]{}
%{
%\begin{minipage}{10cm}
%\color{midnightblue}
%\begin{center}Casa Matriz: Monjitas 527 Piso 15 - Santiago de Chile\\
%Fono (02)24367255 \end{center}
%\end{minipage}
%}
\fancyfoot[C]{}
%Página \thepage\ de \pageref{LastPage}}% Páginas Pares numeradas a la derecha
%\fancyfoot[LO]{\thepage} % Páginas Impares numeradas a la izquierda

\title{C2_Ay5_pauta}

\begin{document}

\begin{center}
Cálculo II\\
Ayudantía Nº5 - Pauta\\
Primer Semestre 2017
\end{center}



\begin{enumerate}
\item Determine si las siguientes afirmaciones son verdaderas o falsas:
\begin{enumerate}[label=\emph{\alph*)}]
\item $\dis \int_0^1 x\, dx \leq \int_0^1 x^2\, dx$ \\

\begin{align*}
 \dis  0 \leq x \leq 1  &\Rightarrow x^2 \leq x\\
 \dis &\Rightarrow \int_0^1 x^2 dx \leq \int_0^1 x dx
\end{align*}

$\therefore$ La afirmación es \underline{Falsa} \dotfill \textbf{(1 punto)}

\item Si $f$ y $g$ son funciones integrables en todo intervalo $[a,b] \subset R$ tales que
\begin{center}
$\dis \int_1^2 f(x)\, dx=-1  \hspace{1cm}  \int_1^5 f(x)\, dx=3 \hspace{1cm} \dis \int_1^5 g(x)\, dx=1 $
\end{center}
Entonces $\dis \int_2^5 f(x)\, dx=4$ \hspace{2mm} y \hspace{2mm} $\dis \int_1^5 g(x)\left(\int_1^5 f(u)du\right)  \, dx=3$
\end{enumerate}
\vspace{2mm}

\begin{align*}
 \dis  \int_1^5 f(x) dx  &= \int_1^2 f(x) dx + \int_2^5 f(x) dx\\
 \dis  3  &= -1 + \int_2^5 f(x) dx\\
 \dis  \int_2^5 f(x) dx  &= 4\\
\end{align*}
$\therefore$ Esta afirmación es \underline{Verdadera} \dotfill \textbf{(1 punto)}

\begin{align*}
 \dis  \int_1^5 g(x) (\int_1^5 f(u) du) dx  &= \int_1^5 g(x) (3) dx\\
 \dis   &=  (3) * 1\\
 \dis   &= 3\\
\end{align*}
$\therefore$ Esta afirmación es \underline{Verdadera} \dotfill \textbf{(1 punto)}

\newpage







\item Calcular: \\
\[ \lim_{x \to 0} \frac{\ln(1+x^2)+\dis \int_1^{\cos x} e^{-t^2}dt}{xe^{x^3}+x^2-\dis \int_0^{x}e^{t^3}dt}\]
\vspace{2mm}
Aplicando L'Hopital

\begin{align*}
 \dis  L  &= \lim_{x \to 0} \frac{\frac{2x}{1+x^2}-\sin(x)e^{-\cos(x)}}{e^{x^3}+3x^3e^{x^3} +2x-e^{x^3}}\\
  \dis    &= \lim_{x \to 0} \frac{\frac{2x}{1+x^2}-\sin(x)e^{-\cos(x)}}{3x^3e^{x^3} +2x} = \frac{0}{0}
\end{align*}
\dotfill \textbf{(2 puntos)}

Nuevamente aplicando L'Hopital
\begin{align*}
 \dis  L  &= \lim_{x \to 0} \frac{\frac{2-2x^2}{(1+x^2)^2}-[\cos(x)e^{-\cos(x)}-\sin(x)e^{-\cos(x)}2\cos(x\sin(x))]}{9x^2e^{x^3} + 9x^5e^{x^3} +2}\\
  \dis    &= \frac{2-\frac{1}{e}}{2}\\
  \dis    &= 1- \frac{1}{2e}\\
\end{align*}
\dotfill \textbf{(1 punto)}







\newpage
\item Sea $F$ la función definida por $\dis F(x)= \int_{\frac{1}{2}}^{x}f(u)du$, donde $\dis f(u)= \int_0^u \frac{1}{1+t^2}dt$. Determine $F'(1)$.\\
\vspace{2mm}

Sabiedo que:
\begin{align*}
 \dis  f(u)&= \int_0^u \frac{1}{1+t^2}dt\\
  \dis   &= \left.\arctan(t)\right|_0^u\\
  \dis   &= \arctan(u)
\end{align*}
\dotfill \textbf{(1 punto)}

Luego

\begin{align*}
 \dis  F'(x)&= f(x) = \arctan(x)\\
  \dis F'(1)  &= \arctan(1) = \frac{\pi}{4}
\end{align*}
\dotfill \textbf{(1 punto)}

\newpage







\item Sea $\dis f(x)= \int_{0}^{g(x)} \frac{dt}{\sqrt{1+t^3}}$, donde $\dis g(x)= \int_{0}^{\cos(x)} \left(1+ \sin (t^2)\right)dt$. Calcule $\dis f'\left(\frac{\pi}{2}\right)$.\\
\vspace{4mm}

\begin{align*}
 \dis  f'(x)&= \frac{1}{\sqrt{1+g^3(x)}}g'(x)\\
  \dis   g'(x)&= 1+\sin(\cos^2(x))(-\sin(x))
\end{align*}
\dotfill \textbf{(2 puntos)}

Como estamos buscando $\dis f'\left(\frac{\pi}{2}\right)$, al evaluar g(x) en el punto $\frac{\pi}{2}$ los limites de integración quedarían de la siguiente manera:\\

\begin{align*}
 \dis  g(x) &= \int_{0}^{\cos(\frac{\pi}{2})} \left(1+ \sin (t^2)\right)dt\\
 \dis   &= \int_{0}^{0} \left(1+ \sin (t^2)\right)dt\\
  \dis   &= 0\\
\end{align*}
\dotfill \textbf{(1 punto)}

Luego

\begin{align*}
 \dis  f'(x)&= \frac{1}{\sqrt{1+g^3(x)}}g'(x)\\
  \dis  f'(x)&= \frac{1}{\sqrt{1+g^3(x)}}(1+\sin(\cos^2(x))(-\sin(x)))\\
    \dis  f'(\frac{\pi}{2})&= \frac{1}{\sqrt{1+g^3(\frac{\pi}{2})}}(1+\sin(\cos^2(\frac{\pi}{2}))(-\sin(\frac{\pi}{2})))\\
    \dis  f'(\frac{\pi}{2})&= \frac{1}{\sqrt{1}}(1+\sin(0)(-1) = -1\\
\end{align*}
\dotfill \textbf{(1 punto)}

\newpage





\item 
\begin{enumerate}[label=\emph{\alph*)}]
\item Sean $\dis \int_1^{11} f(x)\, dx = 15$ y $\dis \int_1^{5} f(x)\, dx = 8$. Calcule $\dis \int_2^{4} f(3x-1)\, dx$\\

Utilizando la sustitución $u = 3x - 1$ y $du = 3 dx$\\
Si $x=2 \Rightarrow u =5$\\
Si $x=4 \Rightarrow u =11$

Luego:
\begin{align*}
 \dis   \int_2^{4} f(3x-1)\, dx &=  \frac{1}{3}\int_5^{11} f(u)\,du\\
 \dis   &= \frac{1}{3}\int_5^{11} f(u)\, du\\
  \dis   &= \frac{1}{3}(\int_1^{11} f(u)\, du-\int_1^{5} f(u)\, du)\\
  \dis   &= \frac{1}{3}(15-8) = \frac{7}{3}\\
\end{align*}
\dotfill \textbf{(1 punto)}

\item Calcule $\dis \int_{-1}^{1} \sqrt{\abs{x}-x} \, dx$\\

\begin{align*}
 \dis  \int_{-1}^{1} \sqrt{\abs{x}-x} \, dx &=  \int_{-1}^{0} \sqrt{-x-x} \, dx + \int_{0}^{1} \sqrt{x-x} \, dx\\
 \dis   &= \int_{-1}^{0} \sqrt{-2x} \, dx + \int_{0}^{1} \sqrt{0} \, dx\\
  \dis   &= \int_{-1}^{0} \sqrt{-2x} \, dx 
\end{align*}
\dotfill \textbf{(1 punto)}

Utilizando la sustitución $u=-2x$ y $du=-2dx$


\begin{align*}
 \dis  \int_{-1}^{0} \sqrt{-2x} \, dx &= -\frac{1}{2} \int_{2}^{0} \sqrt{u} \, du\\
 \dis   &= \frac{1}{2} \int_{0}^{2} \sqrt{u} \, du\\
  \dis   &= \frac{1}{2} \frac{2}{3} \left.u^{\frac{3}{2}}\right|_0^2\\
  \dis   &= \frac{1}{3} 2^{\frac{3}{2}} = \frac{2 \sqrt{2}}{3}\\
\end{align*}
\dotfill \textbf{(1 punto)}
\end{enumerate}




\end{enumerate}




\vspace{10mm}




\end{document}
